% !TEX TS-program = xelatex
\documentclass[11pt,a4paper]{article}

\usepackage[a4paper,left=1in,right=1in,top=0.80in,bottom=0.80in]{geometry}
\usepackage{fontspec}
\usepackage{xcolor}
\usepackage{ragged2e}

\setmainfont{Liberation Sans}
\newfontfamily\lettertitlefont{Liberation Serif}

\definecolor{lettertitle}{HTML}{283A50}
\definecolor{datechip}{HTML}{E6E8EB}

\pagestyle{empty}
\setlength{\parindent}{0pt}
\setlength{\parskip}{0pt}
\linespread{1.08}

\begin{document}

{\lettertitlefont\bfseries\fontsize{26pt}{32pt}\selectfont
\color{lettertitle}Letter of Intent\par}

\vspace{16pt}
{\setlength{\fboxsep}{3pt}
\colorbox{datechip}{August 22, 2026}\par}

\vspace{18pt}
Subject: Expression of Interest in Full-Time Research Position within SVD

\vspace{14pt}
Dear Sir/Madam,

\vspace{12pt}
I am writing to express my interest in continuing with DLR's Simulation and
Virtual Design (SVD) department in a full-time position after completing my
Master's degree in Artificial Intelligence. During my work as a Research
Assistant at DLR, I have contributed to surrogate modeling, research-software
integration, and optimization projects, with work on a complex energy
system consisting of multiple interacting components.

\vspace{10pt}
Working with my supervisor, Dr. Loukas Kyriakidis, we investigated hybrid surrogate
modeling methods for systems with limited experimental data. This work addressed
a central data challenge: simulation data is inexpensive to generate but may not
fully represent real-world behavior, whereas accurate experimental data is costly
and scarce. We therefore implemented and evaluated approaches that combine large
simulation datasets with small amounts of experimental data, including transfer
learning, ensemble methods, and weighted optimization.
We applied these methods to a high-temperature heat pump. The resulting surrogate models
predicted the sink-side outlet temperature and electrical power consumption,
demonstrating the potential of hybrid modeling for data-limited industrial applications.

\vspace{10pt} 
I have also contributed to software integration and automated optimization.
I supported the deployment of middleware for exchanging data with
external project partners and developed APIs for retrieving remote data and
integrating it into RUSH-BI, an optimization framework. These activities strengthened my
ability to connect machine-learning research with reliable and reusable
software systems.

\vspace{10pt}
My Master's thesis, titled ``Reinforcement Learning-Based Optimal Setpoint Determination in Energy Utility Systems,'' investigates whether learned
policies can generate feasible and cost-effective operating setpoints more
quickly than repeatedly solving mathematical optimization problems. It is applied to two energy systems: a
battery energy storage system and a nonlinear industrial steam-generation
system. The policies are evaluated against mathematical optimization based on
operating cost, constraint satisfaction, and computation time.

\vspace{10pt}
In addition to my research experience, I bring more than three years of
experience developing and deploying AI systems using Python, PyTorch, REST
APIs, Docker, and Kubernetes. This combination of machine learning, numerical
optimization, energy-system modeling, and production software engineering
would allow me to contribute across both research development and practical
implementation within SVD.

\vspace{10pt}
Having already worked with the team and gained an understanding of SVD's
research goals, I am motivated to continue contributing further in the department.

\vspace{10pt}
Thank you for considering my application.

\vspace{10pt}
Sincerely

\vspace{8pt}
Ashish Thapa

\end{document}
